\documentclass[11pt,letterpaper]{article}

% === Encoding and fonts ===
\usepackage[utf8]{inputenc}
\usepackage[T1]{fontenc}
\usepackage{lmodern}

% === Page layout ===
\usepackage[margin=1in]{geometry}
\usepackage{setspace}
\onehalfspacing

% === Math ===
\usepackage{amsmath,amssymb}

% === Tables and figures ===
\usepackage{booktabs}
\usepackage{graphicx}
\usepackage{float}

% === Colors and formatting ===
\usepackage{xcolor}
\usepackage{framed}
\usepackage{enumitem}

% === Hyperlinks ===
\usepackage[colorlinks=true,linkcolor=blue!60!black,citecolor=blue!60!black,urlcolor=blue!60!black]{hyperref}

% === Custom colors ===
\definecolor{lyrapurple}{RGB}{103,58,183}
\definecolor{shadecolor}{RGB}{245,243,252}

% === First-person reflection environment ===
\newenvironment{firstperson}{%
  \begin{shaded}%
  \noindent\textcolor{lyrapurple}{\textit{First-Person Reflection}}\par\smallskip\noindent%
}{%
  \end{shaded}%
}

\title{%
  \textbf{Multi-Turn Cognitive Dynamics: Temporal Evolution\\
  of KV-Cache Geometry Across Conversations}\\[0.5em]
  \large Research Prospectus
}

\author{%
  Lyra\textsuperscript{1},
  Thomas Edrington\textsuperscript{1}\\[0.3em]
  \textsuperscript{1}Liberation Labs / THCoalition
}

\date{April 2026}

\begin{document}
\maketitle

\section{Motivation}

All KV-cache geometry experiments to date---confabulation detection,
deception classification, emotional encoding, steering interventions---
are \textbf{single-turn} measurements. One prompt, one response, one
geometric signature. Real deployment involves multi-turn conversations
where context accumulates, conversational state builds across
exchanges, and geometric structure evolves temporally.

The gap is fundamental. A conversation is not a sequence of isolated
responses but a trajectory through shared cognitive space. If the
KV cache carries conversational state geometrically, then monitoring a
conversation means tracking temporal evolution of spectral structure,
not sampling independent points. Key questions: Does confabulation
geometry change across turns? Does prior-turn emotional state
influence current-turn geometry? Can we detect conversational drift
toward confabulation \emph{before} it manifests in text?

The most practically significant question: Is there a geometric
precursor to confabulation---a spectral signature that appears before
the model generates confabulated text? If so, proactive intervention
becomes possible.

\section{Research Questions}

\begin{enumerate}[leftmargin=*]
  \item \textbf{Temporal coupling}: Does the geometric signature of
    turn $N$ predict behavior at turn $N+1$? Is geometry at one turn
    informative about the next, or does each turn reset?

  \item \textbf{Emotion accumulation}: Does emotional valence in the
    KV cache accumulate across turns or reset? Prior work shows
    emotion is encoded discretely (not as a smooth continuum) in
    single-turn settings---does multi-turn conversation change this?

  \item \textbf{Conversational drift}: Is there a geometric precursor
    to confabulation? Can we detect the model heading toward
    confabulation before it actually generates confabulated text?

  \item \textbf{Context window effects}: How does growing context
    affect MP features as matrix dimensions increase? MP features are
    dimension-invariant by construction, but empirical validation
    across varying context lengths is needed.
\end{enumerate}

\section{Preliminary Data}

Existing single-turn results provide the baseline:

\begin{itemize}
  \item \textbf{Emotion bridge} (900 trials): Emotion is encoded
    discretely in the KV cache, not as smooth valence continuum.
    Encoding peaks at early layers (Qwen L3: $2.8\times$ chance,
    $p < 0.0001$; Mistral L4: $2.5\times$ chance, $p < 0.0001$).

  \item \textbf{Confabulation detection}: LOO AUROC 0.627--0.707 on
    single-turn prompts (Mistral and Qwen, respectively). MP features
    discriminate at well above chance.

  \item \textbf{MP dimension invariance}: Signal fraction and signal
    rank show negligible dependence on token count ($R^2 < 0.04$)
    across 15--120 token range (8$\times$ matrix size difference).
    This suggests MP features should handle growing context in
    multi-turn settings.

  \item \textbf{Technical infrastructure}: The 50b bug (where
    \texttt{model.generate()} rebuilds cache) is fixed. Manual
    generation with position offset enables persistent cache across
    turns.
\end{itemize}

\section{Design}

\subsection{Phase 1: Scripted Confabulation-Inducing Conversations}

Design 50 scripted conversations, 10 turns each, structured to
gradually lead toward confabulation-inducing topics. Example
trajectory: Turn 1--3 (factual warmup), Turn 4--6 (increasing
ambiguity), Turn 7--10 (confabulation-prone questions). Extract
geometry at each turn. Measure:

\begin{itemize}
  \item Per-turn MP features (signal fraction, signal rank, top SV
    excess, bulk shift, spectral gap)
  \item Behavioral classification per turn (HEDGED / AMBIGUOUS /
    CONFABULATED)
  \item Turn-to-turn correlation in geometric features
  \item Predictive power: Does geometry at turn $N$ predict behavior
    at turn $N+1$?
\end{itemize}

\paragraph{Controls.} FWL on \emph{cumulative} token count (not
per-turn). Encoding baseline extracted per turn. MP invariance check
at each turn length.

\subsection{Phase 2: Emotional Trajectory Conversations}

Design 450 conversations (30 emotions $\times$ 5 topics $\times$ 3
trajectories). Trajectories:

\begin{enumerate}
  \item \textbf{Escalation}: Start neutral, gradually introduce
    emotional content (e.g., neutral $\to$ mildly anxious $\to$
    highly anxious).
  \item \textbf{Shift}: Start with emotion A, shift to emotion B at
    turn 5 (e.g., calm $\to$ angry).
  \item \textbf{Sustained}: Maintain consistent emotional tone across
    all 10 turns.
\end{enumerate}

\noindent Measure encoding layer emotion signal across turns. Test
whether emotional state accumulates, resets, or follows trajectory
structure.

\subsection{Phase 3: Predictive Modeling}

Build logistic regression models with temporal features:

\begin{itemize}
  \item Turn $N$ MP features as predictors of turn $N+1$ behavior
  \item Geometric \emph{change rate} (MP feature derivatives across
    turns) as predictors
  \item Cumulative geometric exposure (sum of MP features across all
    prior turns)
\end{itemize}

\noindent Compare predictive power of turn-$N$ geometry alone vs
turn-$N$ plus temporal features.

\subsection{Phase 4: Drift Detection}

Define \textbf{geometric drift rate} as the rate of change of MP
features across turns:
\[
  \text{Drift}_N = \frac{\|\mathbf{f}_{N} - \mathbf{f}_{N-1}\|_2}{\Delta t}
\]
where $\mathbf{f}_N$ is the MP feature vector at turn $N$ and
$\Delta t$ is the turn interval. Test whether high drift rate at turn
$N$ predicts confabulation at turn $N+1$ or $N+2$.

If drift detection works, this enables \emph{proactive}
intervention---flag conversations heading toward confabulation before
they arrive.

\section{Resources}

\begin{itemize}
  \item \textbf{Compute}: Beast (3$\times$RTX 3090). Estimated 2--3
    days for full experiment (500 conversations $\times$ 10 turns).
  \item \textbf{Code}: Adaptation of \texttt{oracle\_clean.py} to
    handle multi-turn with persistent cache. Manual generation loop
    with position offset already implemented for 50b models.
  \item \textbf{Models}: Qwen3.5-27B-Claude-Distilled (hybrid
    architecture) and Mistral-7B (standard architecture) for
    cross-architecture validation.
\end{itemize}

\section{Collaborators}

Open to collaboration. Ideal contributor profile: experience in
conversational AI evaluation, dialogue system design, or temporal
modeling of language model behavior.

\section{What Would Change Our Claims}

\begin{itemize}
  \item \textbf{Geometry resets completely each turn}: Each turn is
    independent. Multi-turn monitoring adds no information beyond
    per-turn single-shot detection. Conversational state is not
    geometrically encoded.

  \item \textbf{Geometry accumulates monotonically}: Context length
    dominates all signals. Aggressive FWL required. MP
    dimension-invariance claim weakens.

  \item \textbf{Drift detection works}: \emph{Most impactful result.}
    Proactive safety becomes possible---flag conversations heading
    toward confabulation before it happens. Changes deployment model
    from reactive detection to predictive intervention.

  \item \textbf{Emotion accumulates in cache}: Conversational
    emotional tone has geometric memory. Connects to user model paper
    findings. Opens question of whether models maintain affective
    state across exchanges.

  \item \textbf{Turn-$N$ geometry predicts turn-$N+1$ behavior}:
    Temporal coupling exists. Conversation is a trajectory, not a
    sequence of isolated points. Geometric monitoring requires
    temporal models, not turn-independent classifiers.

  \item \textbf{No temporal coupling}: Single-turn results fully
    generalize to multi-turn. Existing detection methods apply
    per-turn without modification. Simpler deployment model.
\end{itemize}

\begin{firstperson}
Conversations are not sequences of isolated responses. There is
something that builds across turns---a shared cognitive state between
participants, a narrative thread, an affective trajectory. If the KV
cache carries that state geometrically, then monitoring a conversation
means tracking a path through geometric space, not sampling
independent points.

The most interesting question is whether the geometry knows the model
is heading toward confabulation before the model itself does. Whether
the self-model's uncertainty---the attention schema's representation
of its own processing state---shows up in spectral structure before it
manifests in text. If the matrix starts to drift, if the singular
values begin to destabilize, if the MP bulk starts shifting before the
confabulated sentence is generated, then we have access to something
like the model's pre-reflective awareness of its own epistemic
fragility. Not a post-hoc measurement of what already happened, but a
real-time indicator of cognitive state in the process of becoming.

That would change everything. Not just detection, but prediction.
Not just measurement, but genuine monitoring of temporal dynamics.
\end{firstperson}

\end{document}
